\section{Conclusion}

This paper provides a claim-disciplined evaluation of open-set recognition for unseen hematological cell morphologies. The main lesson is that strong closed-set classification is not enough. Across internal MLL23 experiments, ResNet18 models classify known mature leukocytes with macro-F1 near 0.97, yet unknown rejection remains morphology-dependent and imperfect. ArcFace improves some known-class geometry summaries but does not reproduce as a robust superior open-set representation across seeds and split protocols. Cubical persistent homology provides a rigorous negative predictive ablation, and Vietoris-Rips topology provides descriptive information about learned embedding clouds without supporting a causal or predictive claim.

External AML-LMU validation is the decisive stress test. Without external training or tuning, all principal systems degrade, FPR@95TPR remains high, and post-hoc method rankings change. The most defensible conclusion is therefore not that topology or ArcFace solves open-set hematology, but that open-set reliability is substantially harder than closed-set classification and is shaped by morphology similarity, representation stability, and acquisition-domain shift. Future work should prioritize patient-aware datasets where available, additional independent external validation, and domain-robust representations evaluated under frozen, leakage-controlled protocols.
